\pdfoutput=1
\documentclass[11pt]{amsart}
\usepackage{amsmath,amssymb,amsthm}
\usepackage[margin=1.2in]{geometry}
\usepackage[colorlinks=true,linkcolor=blue,citecolor=blue,urlcolor=blue]{hyperref}
\hypersetup{pdfauthor={Bernd Johannes Wuebben},
  pdftitle={Link periods and deformations of Calabi-Yau threefolds with toric singularities}}
\setlength{\emergencystretch}{2em}
\newtheorem{theorem}{Theorem}[section]
\newtheorem{proposition}[theorem]{Proposition}
\newtheorem{lemma}[theorem]{Lemma}
\newtheorem{corollary}[theorem]{Corollary}
\theoremstyle{definition}
\newtheorem{example}[theorem]{Example}
\theoremstyle{remark}
\newtheorem{remark}[theorem]{Remark}
\newcommand{\C}{\mathbb C}
\newcommand{\Q}{\mathbb Q}
\newcommand{\Z}{\mathbb Z}
\newcommand{\OO}{\mathcal O}
\DeclareMathOperator{\Def}{Def}
\DeclareMathOperator{\Ext}{Ext}
\DeclareMathOperator{\Spec}{Spec}
\DeclareMathOperator{\Gr}{Gr}
\DeclareMathOperator{\PD}{PD}
\DeclareMathOperator{\im}{im}
\title[Link periods and Calabi--Yau deformations]{Link periods and deformations of Calabi--Yau threefolds with toric singularities}
\author{Bernd Johannes Wuebben}
\date{November 26, 2026\newline This version: November 26, 2026}
\subjclass[2020]{14B07 (primary); 14J32, 14M25, 32G20 (secondary)}
\keywords{Toric singularities, link periods, Calabi--Yau threefolds, semiuniversal deformations, Minkowski decompositions}
\begin{document}
\begin{abstract}
The periods of a relative volume form on the link of an isolated Gorenstein
toric threefold singularity embed its semiuniversal deformation germ in
the third cohomology of the link. For the homogeneous toric form this
embedding is linear in Altmann's parameters. Multiplication by a
holomorphic unit changes the embedding by an automorphism of its source,
even over nonreduced bases. The automorphism comes from a common change
of the toric height coordinate at all edges of the defining polygon.
Let $X$ be a connected normal projective threefold with trivial
dualizing sheaf, vanishing irregularity and only such singularities.
In the product of the local period images, impose the linear condition
that the Poincar\'e dual link classes sum to zero in the second homology
of a projective log resolution. Using a local-to-global lifting theorem
for small Artin extensions, we identify the semiuniversal germ of $X$ with
the product of this scheme-theoretic linear section and a smooth germ
of dimension $h^1(X,\Theta_X)$. We construct a projective example in
which a first-order deformation extends to second order and is
obstructed at third order.
\end{abstract}
\maketitle
\enlargethispage{2pt}

\section{Introduction}

At an ordinary double point of a Calabi--Yau threefold, the period of a
volume form over the vanishing three-sphere is a local smoothing
parameter. When several nodes lie on a compact threefold, their
parameters satisfy the homological relations among the exceptional
curves of a small resolution \cite{Fri,LLW}. For a general isolated
Gorenstein toric singularity in dimension three, the local deformation
space can itself be reducible or nonreduced. Altmann describes its
equations in terms of a lattice polygon \cite{Alt}. To use these
equations on a compact threefold one must relate the polygon parameters
at different singularities to the periods of a single global form.
Even locally, the restriction of that form can differ from the toric
one by a unit depending on the point of the threefold as well as on
the deformation parameters.

\subsection*{The local period embedding}
Let $Q$ be a lattice polygon with primitive edges, and let $X_Q$ be
the affine toric threefold defined by the cone over $Q$ at height one.
These are the local models for isolated Gorenstein toric singularities
in dimension three. Write $d_1,\ldots,d_n$ for the oriented edge
vectors in cyclic order. Altmann's semiuniversal base $B_Q$ is the
analytic germ at zero given by
\begin{equation}\label{eq:moments}
 \sum_{i=1}^n d_i t_i^k=0 \quad(k\geq 1),\qquad
 (t_1,\ldots,t_n)\quad\text{modulo common translation}.
\end{equation}
We take the slice $t_1=0$. The equations of degree one embed $B_Q$ in
its tangent space $W_Q$, of dimension $n-3$. For a lattice Minkowski
decomposition of $Q$, the edges belonging to the same summand have
equal parameters. The reduced components arise from maximal such
decompositions \cite[\S\S2,8]{Alt}. The higher powers in
\eqref{eq:moments} also determine deformations over Artin algebras
that cannot be recovered from these components.

Let $L$ be the oriented link of the vertex. On a deformation, the
restriction of a relative dualizing generator to a neighbourhood of
$L$ gives a class in $H^3(L,\C)$, with the marking induced by the
central fibre. Over a local Artin $\C$-algebra $A$ this is a class in
$H^3(L,\C)\otimes A$. We write $P_g$ for the period map of
$g\Omega_Q$, where $\Omega_Q$ is the homogeneous toric volume form
on Altmann's family and $g$ is a holomorphic unit. We allow $g$ to
depend on a smooth parameter germ $D$, to accommodate the restriction
of a volume form from a larger family.

\begin{theorem}\label{thm:local}
There is a linear isomorphism $\ell_Q:W_Q\to H^3(L,\C)$ whose
restriction to $B_Q$ is the period map of $\Omega_Q$.
For every $g$ as above there is an analytic automorphism
$\varphi_g$ of $B_Q\times D$ over $D$ such that
\[
 P_g(t,u)=\ell_Q\bigl(\operatorname{pr}_1\varphi_g(t,u)\bigr).
\]
This equality gives the marked period class on every Artin pullback.
The automorphism preserves the ideal of each diagonal $t_i=t_j$.
\end{theorem}

The distinction between a change of parameters in the ambient vector
space and an automorphism of $B_Q$ already occurs for a polygon with
six edges. Its base is a point of length ten with nonzero cubic
products. A quadratic change with identity differential takes one of
the defining equations to a nonzero cubic in its local ring
(Section~\ref{sec:examples}). For a volume form the change has a
special shape. The quotient by the horizontal two-dimensional torus
has a height coordinate $z$. Averaging the unit $g$ over the compact
torus gives a function of $z$ and the parameters, with the same effect
on periods as $g$. Let $H$ be its primitive in $z$, vanishing at zero.
The change on the base is $t_i\mapsto H(t_i,t,u)$ for every edge.
The quotients
\[
  \frac{H(z,t,u)-H(t_i,t,u)}{z-t_i}
\]
are units. The moment equations imply all the multiplicative relations
these units must satisfy to lift $z\mapsto H(z,t,u)$ through the
binomial equations of the total family. This lift has the required
Jacobian on the volume form (Section~\ref{sec:local}). Homogeneity
makes the periods of $\Omega_Q$ linear. Identifying their differential
with the logarithmic Hodge comparison on the link proves that these
linear forms give coordinates (Section~\ref{sec:hodge}).

\subsection*{The compact threefold}
Let $X$ be a connected normal projective complex threefold with
$\omega_X\simeq\OO_X$, $H^1(X,\OO_X)=0$ and isolated singularities
analytically isomorphic to germs $(X_{Q_p},0)$ as above. Write
$\Def(X)$ for its semiuniversal analytic germ, with the central
fibre identified with $X$. Fix local toric
identifications and a projective log resolution $Y\to X$ that is an
isomorphism on the regular locus and has reduced exceptional divisor
with simple normal crossings. Realize each link $L_p$ as the boundary
of a small neighbourhood of the exceptional fibre. Set
\begin{equation}\label{eq:K}
 V_p=H_2(L_p,\C),\qquad V=\bigoplus_p V_p,\qquad
 K=\ker\!\left(V\xrightarrow{\,\sum(\iota_p)_*\,}H_2(Y,\C)\right),
\end{equation}
where $\iota_p:L_p\hookrightarrow Y$ is the inclusion.
Link Poincar\'e duality carries $\ell_{Q_p}(B_{Q_p})$ to a closed
analytic subspace $G_p\subset V_p$. The homogeneous equations of
$G_p$ make it independent of a scalar normalization of $\Omega_{Q_p}$.

\begin{samepage}
\begin{theorem}\label{thm:global}
Suppose that a deformation of $X$ lifts across a small extension of
local Artin $\C$-algebras whenever all its singular germs admit lifts
across that extension. With $h=h^1(X,\Theta_X)$, there is an
isomorphism of analytic germs
\begin{equation}\label{eq:compact}
 (\Def(X),0)\simeq
 \left(\C^h\times\Bigl(K\times_V\prod_p G_p\Bigr),0\right).
\end{equation}
The projection to $V$ consists of the Poincar\'e duals of the link
periods of one global relative volume form. The fibre product is taken
with the scheme structures defined by \eqref{eq:moments}.
\end{theorem}
\end{samepage}

The lifting hypothesis is the local detection theorem of
\cite[Lemma~3.1]{Wue} for projective threefolds with isolated
Gorenstein canonical singularities. It applies here because the
Gorenstein toric germs are canonical. A small extension has kernel
annihilated by the maximal ideal. Existence of a global lift across
each such extension suffices in Theorem~\ref{thm:global}.
To see why, restrict the volume form to the complement of the
singularities. Its boundary classes lie in $K$ with coefficients in
every Artin algebra. The logarithmic description of the derivative
identifies $K$ with the image of tangent localization. Thus, once a
lift exists, a first-order deformation with values in the kernel of
the extension supplies exactly the missing change in its periods.
This gives equality of the analytic ideals in \eqref{eq:compact}
(Section~\ref{sec:compact}).

In particular one can read the equations on $K$ by restricting the
polygon equations, in the linear coordinates $\PD\circ\ell_{Q_p}$.
If $K=0$, only the smooth factor remains. If $K=V$, all the local
deformation parameters occur independently. For intermediate $K$,
the multiplication of nilpotents is determined by restriction of the
ideals. The construction requires a projective log resolution, which
may have nonzero discrepancies.

Gross constructed a projective Calabi--Yau threefold with two cones over
$\mathbb F_1$ and a tangent direction obstructed at second order
\cite[Example~2.19]{Gross}. For an example with a cubic obstruction,
consider a general anticanonical hypersurface in the product of a
toric threefold with $\mathbb P^1$. For the threefold constructed in
Section~\ref{sec:compact-example}, the hypersurface has two singular
points with the same hexagonal polygon. Projection to the toric factor
identifies their links. Under this identification the kernel $K$ is
the antidiagonal, so the compact germ contains one copy of the full
local deformation space as a factor. Its ring, up to smooth parameters, is
\[
 \C\{u,v,s\}/(us,vs,s^3).
\]
The $s$-direction extends to order two, and the equation $s^3=0$
obstructs every extension of that direction to order three.

In general the first obstruction to extending a tangent vector occurs
at the degree of the first polygon moment that is nonzero on its link
periods (Proposition~\ref{prop:jet}). For every $r\geq 2$ we give an
isolated toric singularity with deformation base supported at a point
and a tangent vector whose first obstruction occurs at order $r$
(Proposition~\ref{prop:orders}).

Gross's local-to-global comparison for canonical Calabi--Yau threefolds
gives equations for the reduced support when tangent localization has
positive corank \cite[Theorems~1.16 and~2.2]{Gross}. In the formulation
of Friedman \cite[Theorem~3.8 and Corollary~3.11]{Fried}, generalized
$T^1$-lifting gives this comparison and a product decomposition when
tangent localization is surjective. For nodes and the
specified degree-five and degree-six del Pezzo cones, the description
by marked link periods in \cite[Theorem~1.1]{Wue} also identifies
particular discriminant loci. Here the polygons are arbitrary,
including those with an intrinsically nonreduced local base. Gross's
volume form on the family associated to a Minkowski decomposition
\cite[Proposition~3.2]{GrossSL} gives the toric expression on such a
component. The height change in Theorem~\ref{thm:local} acts on the
whole semiuniversal family, which permits a common global volume form
even in these nonreduced directions.

\section{Changes of height on the toric deformation}\label{sec:local}

Let $Q$ be a lattice polygon with primitive edges in the rank-two
lattice $N$. Translate a vertex to the origin. Write $d_1,\ldots,d_n$
for the oriented edges in cyclic order and $M=\operatorname{Hom}(N,\Z)$.
For Altmann's family $\pi_Q:\mathcal X_Q\to B_Q$ we take the section
$t_1=0$ of simultaneous translation of the edge parameters. Thus
\begin{equation}\label{loc:moments}
 \OO_{B_Q,0}=
 \C\{t_2,\ldots,t_n\}\big/
 \left(\sum_{i=1}^n d_i t_i^k\ :\ k\geq 1\right).
\end{equation}
The first moment identifies the ambient space with $W_Q$, the tangent
space of $B_Q$. For the remaining moments the scheme structure in
\eqref{loc:moments} will matter throughout. The variable $z$ along the
orbits of simultaneous translation is a function on the total space.
We use the convention $\tau_i=z-t_i$, so that $\tau_1=z$.

\subsection{The lift to the total space}
Put $S_\Z=\ker(\Z^n\to N)$, where the $i$th basis vector maps to
$d_i$, and put $\Lambda=\{r\in\Z^n:r|_{S_\Z}=0\}$, using the
coordinate pairing on $\Z^n$. Then $S_\Z^\vee=\Z^n/\Lambda$.
We recall the monomials in Altmann's presentation
\cite[\S\S4.4--5.3]{Alt}. For $c\in M$ choose a vertex of $Q$ at
which $c$ attains its minimum and an edge path from $0$ to that
vertex. If $\lambda_i$ is the signed number of traversals of edge $i$,
the vector with entries $-\lambda_i\langle c,d_i\rangle$ represents
a class $\eta(c)\in S_\Z^\vee$ independent of the choices. A finite
set of monomials $Z_c$ with exponents $(c,\eta(c))$, together with the
$\tau_i$ with exponents $(0,\bar e_i)$, generates the semigroup algebra
used in the construction. Its relations are the equalities of monomials
having the same exponent in $M\oplus S_\Z^\vee$. Base change along the
untranslated moment scheme $B_Q\times\C_z$ gives the total space of
$\pi_Q$ near the central fibre \cite[Theorem~5.1 and \S7]{Alt}.
In particular the presentation uses the subsemigroup generated by these
monomials, even when its semigroup algebra fails to be normal.

Let $T=\operatorname{Hom}(M,\C^*)$ act through the first component of
the exponent. It fixes $z$ and the base. On the open set
$F=\prod_i\tau_i\ne0$ the torus characters can be written
\begin{equation}\label{loc:characters}
             \chi^c=Z_c\prod_i\tau_i^{-\eta_i(c)} .
\end{equation}
For an arbitrary $c$, the notation $Z_c$ denotes the monomial of
exponent $(c,\eta(c))$ in the same algebra
\cite[\S4.6 and \S5.3, Remark~(1)]{Alt}. Here $(\eta_i(c))$ is any
integral representative of $\eta(c)$.
After inverting $F$, all the base monomials are invertible and
the localized semigroup splits off $M$. Indeed the weights $c$ of the
generators generate $M$ as a monoid, since any $c\in M$ occurs in
the dual of the cone over $Q$ at sufficiently large height.
This gives the product with $T$ over the open subset of
$B_Q\times\C_z$ where $F\ne0$.

\begin{proposition}\label{loc:lift}
Let $u$ be coordinates on a smooth germ $D$. Let $H(s,t,u)$ be
holomorphic near the origin, with $H(0,t,u)=0$ and $H_s(0,0,0)\ne0$.
The substitutions
\begin{equation}\label{loc:change}
    t_i'=H(t_i,t,u),\qquad z'=H(z,t,u)
\end{equation}
lift to an automorphism $\Psi_H$ of $\mathcal X_Q\times D$ which fixes
the characters \eqref{loc:characters} on their domain. Its base
automorphism $\varphi_H$ fixes $u$ and preserves the ideal of each
diagonal $t_i=t_j$.
\end{proposition}

\begin{proof}
The moment equations have the following consequence for any convergent
series $f(s)$ whose coefficients depend holomorphically on $t,u,z$:
\begin{equation}\label{loc:series}
              \sum_i r_i f(t_i)=0\qquad(r\in\Lambda).
\end{equation}
Indeed $r$ belongs to the complex row space of the edge matrix, so
every positive power of $s$ gives zero by \eqref{loc:moments}. The
constant term gives zero because $\sum_i d_i=0$. The same calculation
holds in each maximal-ideal quotient, in which it requires only
finitely many powers. Separatedness proves the identity in the analytic
ring. Taking $f(s)=H(s,t,u)^k$ for the two edge rows shows that
\eqref{loc:change} defines a map of $B_Q\times D$.

Set
\[
     U_i=\frac{H(z,t,u)-H(t_i,t,u)}{z-t_i}.
\]
Each divided difference extends holomorphically across $z=t_i$ and
has value $H_s(0,0,0)$ at the origin. Choose a logarithm near this
common value. Apply \eqref{loc:series} to the logarithm of the
divided difference with $s$ in place of $t_i$. Exponentiating gives
\begin{equation}\label{loc:units}
               \prod_i U_i^{r_i}=1 \qquad(r\in\Lambda).
\end{equation}
Thus the assignment $\bar e_i\mapsto U_i$ is a homomorphism from the
whole exponent lattice $S_\Z^\vee$ to the units of the analytic
ring. In particular it respects relations in the saturation of the
integer row span. Multiply each monomial of exponent $(c,\eta)$ by
the value of this homomorphism at $\eta$. On generators this reads
\[
        \tau_i\longmapsto\tau_i U_i,\qquad
        Z_c\longmapsto Z_c\prod_i U_i^{\eta_i(c)} .
\]
Both sides of a binomial acquire the same factor. Since
$\tau_iU_i=z'-t_i'$, these assignments descend to the base change
defining $\mathcal X_Q$ and fix the characters \eqref{loc:characters}.

To invert the construction, lift $H$ to the smooth ambient parameter
space, preserving $H(0,t,u)=0$. There the differential of
$t\mapsto t'$ at the origin is $H_s(0,0,0)$ times the identity in the
$t$ directions. Write $t=b(t',u)$ for its inverse, and let
$J(s',t',u)$ be the inverse in $s$ of $H(s,b(t',u),u)$.
Then $t_i=J(t_i',t',u)$. Formula \eqref{loc:series} for the primed
parameters shows that this inverse descends to $B_Q\times D$.
The divided differences of $J$ pull back to $U_i^{-1}$, so its lift
inverts the map on the total space as well. Finally $t_i'-t_j'$ is
$t_i-t_j$ times a unit, including when $i$ or $j$ equals $1$.
\end{proof}

\subsection{The invariant volume form}
Fix an oriented basis $c_1,c_2$ of $M$. On the open torus family take
the relative form
\begin{equation}\label{loc:form}
       \Omega_Q=dz\wedge d\log\chi^{c_1}\wedge d\log\chi^{c_2}.
\end{equation}
All differentials here and below are relative to $B_Q$ (and to $D$
after taking the product). On families from a fixed Minkowski
decomposition this is Gross's volume form
\cite[Proposition~3.2]{GrossSL}.

\begin{lemma}\label{loc:generator}
The form $\Omega_Q$ generates the relative dualizing line of $\pi_Q$
near the vertex. The automorphism of Proposition~\ref{loc:lift}
satisfies $\Psi_H^*\Omega_Q=H_z(z,t,u)\Omega_Q$, where $H_z$ denotes
the derivative in the height variable.
\end{lemma}

\begin{proof}
Take first a base change to a local Artin $\C$-algebra $A$ with maximal
ideal $\mathfrak m$. The facet opens of the central toric cone cover
the complement of the vertex. For the facet spanned by
$(v_i,1),(v_{i+1},1)$, primitivity of $d_i$ extends this pair to a
lattice basis. The corresponding affine chart is
$\mathbb A^2_{x,y}\times\mathbb G_{m,\zeta}$. In the dual coordinates
the height monomial is $xy$ after absorbing a power of $\zeta$ into
$y$. With a character basis of $T$ adapted to this chart the weights
of $x,y,\zeta$ are $(1,0),(-1,0),(0,1)$.

Lift each facet open by localizing at a semi-invariant lifting its
defining monomial. The resulting affine $A$-algebra is flat, and one
can lift $x,y,\zeta$ in their weight spaces. The lift of $\zeta$ is a
unit. The resulting map from $A[x,y,\zeta^{\pm1}]$ to this algebra
is an isomorphism: on each graded piece for the $\mathfrak m$-adic
filtration it is the identity on the central chart tensored with
$\mathfrak m^k/\mathfrak m^{k+1}$. These algebraic charts apply to
every Artin point of the analytic base, since \eqref{loc:moments} and
the original family are algebraic. In the lifted chart the $T$-invariant
height $z$ is a polynomial $f_i(q)$ in $q=xy$, with
$f_i(q)\equiv q\bmod\mathfrak m$. Its derivative is a unit.
Where $F$ is inverted, $q$ and $x$ are units as well, since $F$
reduces to $q^n$. The characters there are monomials in $x,\zeta$
times invariant units, and the exponents form a unimodular matrix.
The invariant factors have logarithmic differentials proportional to
$dq$, since the invariant ring on this open is $A[q,q^{-1}]$.
Consequently \eqref{loc:form} is
$\pm f_i'(xy)\,dx\wedge dy\wedge d\zeta/\zeta$ on the chart. It
extends as a generator on the punctured fibre over $A$.

Flatness and the Gorenstein property of the central cone give a
relative dualizing line near the vertex, compatible with base change
\cite[Tags 0E4C and 0C08]{Stacks}. On the Artin pullback the local
ring at the vertex has depth three, by flatness and the
Cohen--Macaulay property of the central fibre. The ratio of the form
to a local generator and its reciprocal therefore extend uniquely
across that point.

For descent to the analytic family, choose a local algebraic
generator of the relative dualizing line. In its trivialization
\eqref{loc:form}, which is regular on the algebraic localization at
$F$, has expression $a/F^m$ for some $m$. Multiplication by $F$ is
injective on every Artin pullback: this follows on its associated
graded ring from injectivity of multiplication by $z^n$ on the central
cone. The extension just proved thus gives
$a\in(F^m)+\mathfrak m_{B_Q}^k\OO_{\mathcal X_Q}$ for every $k$.
Krull intersection in the quotient by $(F^m)$ gives $a\in(F^m)$.
The resulting regular section is a generator by its central value.
The formula for $\Psi_H^*\Omega_Q$ follows from
\eqref{loc:form} on inverting $F$. It holds on the family since
injectivity on all the Artin pullbacks and separatedness give
injectivity of multiplication by $F$ in the local ring.
\end{proof}

\subsection{Periods on the link}
Let $L$ be the link with the marking of the central fibre. Over a
local Artin algebra $A$, the deformation is smooth over $A$ in an
annular neighbourhood of $L$. Relative analytic coordinates there
identify the relative de Rham complex locally with the ordinary
holomorphic complex of a polydisc tensored with $A$. The inclusion of
the constant sheaf $A$ is consequently a quasi-isomorphism. Taking
inverse images of sheaves on $L$ gives for any relative three-form
$\omega$ a class
\begin{equation}\label{loc:artin}
                  [\omega]_{L,A}\in H^3(L,\C)\otimes_\C A .
\end{equation}
These identifications commute with maps of Artin algebras. Indeed the
inclusion of constants is natural and a finite good cover computes
constant-sheaf cohomology by a complex of $\C$-vector spaces tensored
with $A$. They also commute with isomorphisms of punctured germs,
using the induced map on central cohomology. One may shrink the annulus
or compute on the punctured germ, since inclusion of the link induces
the same cohomology in either case.

Apply \eqref{loc:artin} to $\Omega_Q$ on the successive maximal-ideal
quotients of $\OO_{B_Q,0}$ to obtain a formal map $\widehat P_1$.
The diagonal cocharacter of the exponent lattice gives $\tau_i$ and
$t_i$ weight one and gives $Z_c$ weight $\sum_i\eta_i(c)$. It fixes
$\chi^c$ and hence gives $\Omega_Q$ weight one. This $\C^*$ acts on
the base, whereas $T$ acts fibrewise. On the punctured central cone
the diagonal cocharacter is homotopic to the identity through torus
actions. Naturality of \eqref{loc:artin} yields
$\widehat P_1(\lambda t)=\lambda\widehat P_1(t)$.
In the completion of a standard graded ring defined by the homogeneous
ideal \eqref{loc:moments}, an element of weight one is its linear
term. Thus $\widehat P_1$ is the restriction of a linear map
\begin{equation}\label{loc:linear}
                  \ell_Q:W_Q\longrightarrow H^3(L,\C).
\end{equation}

Let now $g$ be a holomorphic unit on $\mathcal X_Q\times D$ and let
$h$ be its average over the compact subgroup $(S^1)^2\subset T$.
Take an invariant representative neighbourhood for this average.
Each torus element fixes $\Omega_Q$ and acts trivially on central
link cohomology. Thus $[g\Omega_Q]_{L,A}=[h\Omega_Q]_{L,A}$ for each
Artin point. One can integrate the pullbacks on a compact cycle after
taking smooth de Rham representatives with coefficients in the
finite-dimensional space $A$.

In fact $h$ is holomorphic in $z,t,u$. To see this, lift $g$ to an
ambient convergent series in $Z_c,\tau_i,u$ and take its weight-zero
terms. The exponent monoid of such terms has finitely many generators
by Gordan's lemma. The invariant series converges as a series in
the corresponding monomials: if their maximum total degree is $b$,
a monomial of degree $k$ has at least $k/b$ factors in any
factorization. Choosing one factorization for each term and shrinking
the new polydisc gives absolute convergence by the Cauchy bounds on
the original series. By convexity of $\eta$ \cite[\S4.5(iii)]{Alt},
any product of the $Z_c$ of weight zero has exponent $(0,\bar a)$
for a nonnegative integral vector $a$. It equals $\tau^a$ by the
semigroup relations. Substitution in the finite set of generators
expresses $h$ as a convergent series in $\tau_i,u$, since every
nonconstant generator vanishes at the vertex. Moreover $h$ has the
same nonzero value as $g$ there.

\begin{proof}[Proof of Theorem~\ref{thm:local}]
Choose a convergent representative of $h(z,t,u)$ and apply
Proposition~\ref{loc:lift} to
$H(s,t,u)=\int_0^s h(v,t,u)\,dv$.
The central action of $\Psi_H$ on link cohomology is the identity.
Indeed on the central cone, with $a(z)=H(z,0,0)/z$, its substitutions
are $z\mapsto za(z)$ and $Z_c\mapsto Z_c a(z)^{\sum_i\eta_i(c)}$.
Replacing $a(z)$ by $a(rz)$ for $0\leq r\leq1$ gives an isotopy of
germs to diagonal scaling by the constant $a(0)\ne0$. This scaling
lies in the connected grading torus. By Lemma~\ref{loc:generator} and
the functoriality of \eqref{loc:artin}, the marked class of $g\Omega_Q$
at $(t,u)$ is the class of $\Omega_Q$ at $(t',u)$. On every Artin
point it is therefore $\ell_Q(t')$, where $t'_i=H(t_i,t,u)$.
This expression supplies its convergence. On a disc the usual periods
over transported cycles in a smooth collar are holomorphic functions
whose Taylor coefficients are their relative de Rham classes on the
Artin quotients, so it also gives the usual marked periods there.
Proposition~\ref{prop:variation} identifies \eqref{loc:linear} with
the infinitesimal period isomorphism. The period map consequently
identifies $B_Q$ with its linear image after the base automorphism
$\varphi_H$, uniformly over $D$. The assertion on the marked diagonal
ideals follows from Proposition~\ref{loc:lift}.
\end{proof}

\section{The differential of the link period}\label{sec:hodge}

We first compute the coefficient of $\varepsilon$ in a period over the
dual numbers. Its component of Hodge degree two is contraction by the
Kodaira--Spencer class. Logarithmic forms on a resolution identify
this component in the Hodge
filtration of the link. To compare it with the period itself we construct
a holomorphic primitive for every holomorphic three-form on the punctured
toric germ.

\subsection{Logarithmic forms at a toric point}
Let $(S,p)$ be an isolated Gorenstein toric threefold singularity. Take
a sufficiently small contractible Stein representative and write
$S^\circ=S\setminus\{p\}$. On a good resolution $M\to S$, let $D$ be
the reduced exceptional divisor. We take representatives for which $M$
retracts onto $D$ and $M\setminus D=S^\circ$ has the homotopy type of
the link $L$. Put $\mathcal F^a=\Omega_M^a(\log D)$. Unless indicated
otherwise, topological cohomology has complex coefficients.

Steenbrink's toric vanishing theorem gives
\begin{equation}\label{hg:vanishing}
 H^q(M,\mathcal F^a(-D))=0\qquad (q>0)
\end{equation}
on any good resolution \cite[\S2, Theorem~2]{Ste}. Local duality
identifies the dimension of this group with that of
$H_D^{3-q}(M,\mathcal F^{3-a})$ \cite[Definition~1.4]{FL}.
In particular $H^1_D(M,\mathcal F^2)=H^2_D(M,\mathcal F^2)=0$.
The exact sequence for cohomology with support and the restriction
sequence for $D$ now give isomorphisms
\begin{equation}\label{hg:restrictions}
 H^1(S^\circ,\Omega^2_{S^\circ})
 \xleftarrow{\ \sim\ } H^1(M,\mathcal F^2)
 \xrightarrow{\ \sim\ } H^1(D,\mathcal F^2|_D)
 =\Gr_F^2 H^3(L).
\end{equation}
The last equality comes from the logarithmic complex
$\mathcal F^\bullet|_D$ for the link, whose Hodge spectral sequence
degenerates at $E_1$ \cite[\S3]{Ste}.

In fact $H^3(L)$ is pure of type $(2,2)$. Toric singularities are
rational, so $\Gr_F^0H^2(L)=\Gr_F^2H^2(L)=0$ by
\cite[\S3, Lemma~2]{Ste}. The logarithmic complex has no other
possible Hodge degree in cohomological degree two except one.
Hodge symmetry on each weight quotient therefore leaves only type
$(1,1)$ in $H^2(L)$. Link duality identifies $H^3(L,\Q)$ with
$H^2(L,\Q)^\vee(-3)$ as a mixed Hodge structure \cite[\S3]{Ste}.
In particular $\Gr_F^2H^3(L)=H^3(L)$ canonically.

Choose a generator $\omega$ of the dualizing sheaf on $S$. The
depth-three comparison identifies the first-order deformation space
$T^1_{S,p}$ with $H^1(S^\circ,\Theta_{S^\circ})$
\cite[Lemma~1.17 and Remark~4.2]{FL}. Here the usual cotangent-complex
tangent group agrees with $\Ext^1(\Omega^1_S,\OO_S)$ because the negative
cohomology sheaves of the cotangent complex have finite support and
grade three in $\OO_S$. Contraction with $\omega$ on
$S^\circ$, followed by \eqref{hg:restrictions}, defines an isomorphism
\begin{equation}\label{hg:lambda}
 \lambda_\omega:T^1_{S,p}\longrightarrow H^3(L).
\end{equation}

\begin{proposition}\label{prop:variation}
Let $A=\C[\varepsilon]/(\varepsilon^2)$. On a deformation of $(S,p)$
over $A$ with Kodaira--Spencer class $v$, every relative volume form
lifting $\omega$ has marked link class
$\varepsilon\lambda_\omega(v)$. In particular the differential of the
period of the homogeneous toric generator is an isomorphism from
$W_Q$ to $H^3(L)$.
\end{proposition}

\begin{proof}
We shall use the exactness of holomorphic three-forms on $S^\circ$
after shrinking $S$. Write such a form as $f\omega_0$, where $\omega_0$
is the toric generator and $f$ extends across $p$ by normality.
An integral cocharacter in the interior of the cone contracts $S$ to
$p$. Its Euler field $e$ has positive weights on the nonconstant
semigroup generators, and $\omega_0$ has positive weight $a$, the
pairing of the cocharacter with the Gorenstein degree. On a small
invariant polydisc for the contracting real scalars $0<s\leq1$, set
\[
 F(x)=\int_0^1 s^{a-1} f(s\cdot x)\,ds .
\]
In homogeneous ambient coordinates the integral divides the coefficient
of each monomial of weight $m$ by $a+m$, so $F$ is holomorphic. Thus
$(e+a)F=f$ and Cartan's formula gives
$d\iota_e(F\omega_0)=f\omega_0$. In particular choose
$\Gamma$ on $S^\circ$ with $d\Gamma=\omega$.

Trivialize the first-order deformation on a Stein cover of $S^\circ$.
Use the convention that transport of forms from chart $j$ to chart $i$
is $1-\varepsilon\mathcal L_{\kappa_{ij}}$ for the
Kodaira--Spencer cocycle $\kappa$. Write the lifted form in chart $i$
as $\omega+\varepsilon\eta_i$ and put
$\beta_{ij}=\iota_{\kappa_{ij}}\omega$. The gluing identities give
$\delta\eta=d\beta$ and $\delta\beta=0$. With total differential
$\delta+(-1)^q d$ in \v Cech degree $q$, the pair $(\eta,\beta)$
represents the coefficient of $\varepsilon$ in the marked class.
Indeed lift $\Gamma$ by its central expression on each chart.
Subtracting the total coboundary of this lift from the relative volume
form leaves $\varepsilon(\eta,\mathcal L_\kappa\Gamma)$ on the central
complex. The second entry differs from $\beta$ by
$d\iota_\kappa\Gamma$, a total coboundary since
$\delta\iota_\kappa\Gamma=0$. The relative de Rham resolution of the
constant sheaf $A$ identifies multiplication by $\varepsilon$ in this
calculation with the coefficient inclusion $\C\to A$,
$c\mapsto\varepsilon c$.

Lift $[\beta]$ by \eqref{hg:restrictions} and choose a logarithmic
one-cocycle $b$ representing the lift. We claim that $db$ is a
\v Cech coboundary on $M$. By \eqref{hg:vanishing} for $a=3$ the
restriction
$H^1(M,\mathcal F^3)\to H^1(D,\mathcal F^3|_D)$ is an
isomorphism. There the class of $db$ is the first differential of
the link Hodge spectral sequence on $[b|_D]$, hence zero. We may
therefore choose a zero-cochain $c$ of logarithmic three-forms with
$\delta c=db$. The total cocycle $(c,b)$ on $M$ restricts to a
class on the link whose Hodge degree-two component is $[b|_D]$.
As $F^2H^3(L)=H^3(L)$ and $F^3H^3(L)=0$, that component determines
the class and gives $\lambda_\omega(v)$.

The logarithmic comparison identifies restriction of $(c,b)$ to
$S^\circ$ and to the link complex under $H^3(S^\circ)\simeq H^3(L)$
\cite[\S1.1]{FL}. On $S^\circ$ choose a zero-cochain $u$ of
two-forms with $\delta u=\beta-b$, refining the covers as necessary.
The difference of $(\eta,\beta)$ and $(c,b)$, modulo the total
coboundary of $u$, is the global holomorphic three-form
$\eta-c-du$. Its primitive constructed above gives zero in $H^3(L)$.
This proves the formula for every lift of $\omega$. For $S$ the cone
of $Q$, the Kodaira--Spencer isomorphism $W_Q\simeq T^1_{S,p}$
of \cite[\S6.2 and the corollary in \S7.2]{Alt} gives the final assertion.
\end{proof}

\subsection{Restriction from the smooth locus}
Let $X$ be a normal projective threefold with isolated Gorenstein toric
singularities. Choose a global volume form $\omega$ trivializing its
dualizing sheaf. Write
$U=X_{\mathrm{reg}}$. Let $Y\to X$ be a projective good resolution
which is an isomorphism over $U$, with reduced exceptional divisor
$E=\coprod_p E_p$. Orient each $L_p$ as the boundary of a small
resolution neighborhood of $E_p$.

\begin{proposition}\label{prop:tangent}
Composing tangent localization with the local period differentials for
$\omega$ and link Poincar\'e duality gives an exact sequence
\[
 0\longrightarrow H^1(X,\Theta_X)
 \longrightarrow T_0\Def(X)
 \longrightarrow K\longrightarrow 0 .
\]
Under the same duality the image of
$H^3(U)\to\bigoplus_p H^3(L_p)$ is $K$.
\end{proposition}

\begin{proof}
Put $\mathcal F=\Omega_Y^2(\log E)$ and let $Y_p$ be the inverse
image of a small Stein representative at $p$. Excision identifies
$H^i_E(Y,\mathcal F)$ with the direct sum of the corresponding
groups on $Y_p$. Their vanishing for $i=1,2$ proves that
$H^1(Y,\mathcal F)\to H^1(U,\Omega_U^2)$ is an isomorphism.
The natural depth-three comparison
$T_0\Def(X)\simeq H^1(U,\Theta_U)$ intertwines tangent localization
with restriction to the punctured representatives
\cite[Lemma~4.1 and Remark~4.2]{FL}. Under contraction with
$\omega$ and the inverse restriction isomorphisms, localization
therefore becomes the map
\[
 H^1(Y,\mathcal F)\longrightarrow\bigoplus_p H^1(Y_p,\mathcal F|_{Y_p}).
\]

For $j:U\hookrightarrow Y$, restriction of $Rj_*\Q_U$ to $E$ gives the
map $r:H^3(U)\to\bigoplus_p H^3(L_p)$. Its de Rham realization is
restriction of $\Omega_Y^\bullet(\log E)$ to $E$, with the filtration
by form degree \cite[\S1.1]{FL}. In particular $r$ is a morphism of
mixed Hodge structures. By logarithmic Hodge degeneration for the
smooth projective pair $(Y,E)$ \cite[Corollary~3.2.13(ii)]{Del}, the
source of the last display is $\Gr_F^2H^3(U)$. After the local
identifications \eqref{hg:restrictions}, the displayed map is
$\Gr_F^2r$. Its target has only Hodge type $(2,2)$, so strictness
\cite[Theorem~2.3.5(iii)]{Del} gives
$\im(\Gr_F^2r)=\im r$. Proposition~\ref{prop:variation} identifies
this image with that of the differential of the marked periods.

The Mayer--Vietoris sequence for $Y=U\cup\bigcup_pY_p$ contains
\[
 H^3(U)\oplus\bigoplus_p H^3(E_p)
 \longrightarrow \bigoplus_p H^3(L_p)
 \xrightarrow{\partial} H^4(Y).
\]
The arrow $H^3(E_p)\to H^3(L_p)$ respects mixed Hodge structures
\cite[(1.4)]{FL}. Its source has weights at most three since $E_p$ is
proper, whereas the target is pure of weight four. This arrow
vanishes, giving $\im r=\ker\partial$. The boundary map is the sum
of the Gysin maps of the oriented links, as follows by excision and the
Thom isomorphism on their collars. Poincar\'e duality thus identifies
its kernel with the kernel of $\bigoplus_p H_2(L_p)\to H_2(Y)$.
This proves both image assertions. Finally, the local-to-global
spectral sequence for the cotangent complex identifies the kernel of
tangent localization with $H^1(X,\Theta_X)$. Each local period
differential is an isomorphism by Proposition~\ref{prop:variation},
so this is also the kernel in the asserted sequence.
\end{proof}

\section{The equations on a projective threefold}\label{sec:compact}

Let $B=\Def(X)$ and let $\mathcal X\to B$ be a proper semiuniversal
family. We use the local models of Section~\ref{sec:local} with the
markings fixed in the introduction. A choice of global generator of
$\omega_X$ fixes the scale of the infinitesimal period comparison
at every point.

\subsection{The period map of a global form}
After shrinking $B$, the family is relatively Gorenstein and its
relative dualizing sheaf is a line bundle. Its restriction to $X$
has a generator $\omega$. The vanishing $H^1(X,\omega_X)=0$ permits
a lift of $\omega$ over the full analytic germ. Indeed a finite free
complex for the proper cohomology of the relative dualizing sheaf
computes cohomology after base change. In a minimal such complex the
degree-zero term has rank one and the adjacent terms vanish, by the
cohomology of the central fibre. A generator of the resulting direct
image lifts $\omega$, and its zero locus has closed image disjoint
from the origin. Shrinking once more yields a relative volume form
$\widetilde\omega$. The same calculation for $\OO_{\mathcal X}$
gives $\pi_*\OO_{\mathcal X}=\OO_B$, where $\pi$ is the family map.
This form of the base change argument, including nonreduced $B$,
is also used in \cite[\S2]{Wue}.

For each $p$ choose an analytic classifying map $b_p:B\to B_{Q_p}$
and an identification of the germ of the global family at $p$ with
the pullback of the toric family. Embed $B$ in a smooth germ $D$.
The graph of $b_p$ in $B_{Q_p}\times D$ is closed, so the above
pullback is a closed subgerm of $\mathcal X_{Q_p}\times D$. There
$\widetilde\omega/\Omega_{Q_p}$ is a unit on the pullback. Lift it to
a unit $g_p$ on the ambient family by lifting a representative in the
local ring. Theorem~\ref{thm:local} supplies over $D$ an isomorphism
\[
  B_{Q_p}\times D\longrightarrow G_p\times D
\]
whose first component is the marked period after link duality. Its
restriction to the graph of $b_p$ depends only on the restricted form:
on each Artin quotient it is the de Rham class of
$\widetilde\omega$, and the analytic local ring of $B$ is separated.
Taking all these restrictions gives $P:B\to G=\prod_pG_p$.

For any Artin pullback the smooth locus over the central open set
$U=X_{\rm reg}$ carries the closed form $\widetilde\omega$.
The relative de Rham resolution of the constant coefficient ring
therefore puts its link classes in the image of $H^3(U)\to
\bigoplus_p H^3(L_p)$ tensored with that ring. By
Proposition~\ref{prop:tangent} this image is $K$ after duality.
Since the linear equations of $K$ vanish on every maximal-ideal
quotient of $\OO_{B,0}$, the analytic map $P$ factors through
$K\times_V G$. Its tangent map has image $K$ and kernel
$H^1(X,\Theta_X)$ by the same proposition, now applied to the
central form $\omega$. A different relative generator rescales the
tuple by a single unit on $B$, which preserves both the linear
subspace and the homogeneous local ideals.

\subsection{Lifting with prescribed periods}
\begin{proof}[Proof of Theorem~\ref{thm:global}]
Choose functions $s_1,\ldots,s_h$ on $B$ vanishing at the origin
whose differentials restrict to a basis of the dual of $\ker dP$.
They give a map
\[
 F=(s,P):B\longrightarrow T=\C^h\times (K\times_V G).
\]
Since each $G_p$ has tangent space $V_p$, one has
$T_0T=\C^h\oplus K$ and $dF$ is an isomorphism. In particular $F$
is a closed immersion of germs. Indeed the pullbacks of ambient
coordinates span the cotangent space of $B$. A basis among them is a
coordinate system in a minimal smooth embedding of $B$, by the
analytic inverse function theorem.

Let $A'\to A$ be a small extension with kernel $I$ and let $b\in B(A)$.
Fix $y'\in T(A')$ above $F(b)$. Lift the image of $b$ in the
smooth ambient germ $D$ to $D(A')$. With this auxiliary parameter
fixed, the $p$th component of the prescribed period in $y'$ pulls
back under the inverse local period isomorphism to a point of
$B_{Q_p}(A')$ above $b_p(b)$. This proves local liftability of the
given deformation simultaneously at all of its singular points.
The lifting hypothesis now supplies a global lift. A semiuniversal
family induces a smooth morphism from its formal base to the
deformation functor. Hence this lift is represented by a point
$b'\in B(A')$ reducing to the specified point $b$.

The two homomorphisms $\OO_{T,0}\to A'$ defined by $y'$ and
$F(b')$ have the same reduction to $A$. Their difference is a
derivation at the closed point with values in $I$, because
$\mathfrak m_{A'} I=0$. Let $v\in T_0B\otimes I$ be the inverse
image of this difference under $dF\otimes I$. Adding $v$ to the
homomorphism defining $b'$ again gives a ring homomorphism, and
its image under $F$ is $y'$. This proves formal smoothness of the
closed immersion. Applying the lifting to the successive
maximal-ideal quotients of $\OO_{T,0}$ gives a left inverse to the
completed quotient homomorphism, hence its injectivity.
Faithful flatness of completion gives the asserted analytic
isomorphism.
\end{proof}

\begin{remark}
The lifting argument also describes what follows from local
liftability alone. Suppose $B$ classifies a semiuniversal deformation
with isolated singularities and the lifting hypothesis of
Theorem~\ref{thm:global} holds. Write $D_0$ for the product of its
local hulls, $\lambda:T_0B\to T_0D_0$ for tangent localization, and
$h_0,r$ for the dimensions of its kernel and cokernel. In a minimal
smooth embedding of $D_0$ with ideal $I_0$, the global base admits
a presentation
\[
 \OO_{B,0}\simeq \C\{s_1,\dots,s_{h_0},t_1,\dots,t_{\dim T_0D_0}\}
             /(I_0,f_1,\dots,f_r)
\]
with the $df_i(0)$ linearly independent. To see this, embed $B$ in
$\C^{h_0}\times D_0$ by the local classifying maps and functions
dual to the kernel as above. Choose elements of its ideal whose
linear parts span the conormal space in the smooth ambient germ.
Their zero locus in $\C^{h_0}\times D_0$ has the same tangent space
as $B$ and any point of this locus gives local lifts via its
projection to $D_0$. The same small-extension argument proves
equality of the germs. For the toric points, the periods give the
$f_i$ as linear equations on $V$, independent of the smooth
coordinates.
\end{remark}

\section{The node and two nonreduced deformation spaces}\label{sec:examples}

\subsection*{The period at a node}
Write the deformation of the ordinary double point as
\[
 x_1y_1=z,\qquad x_2y_2=z-s.
\]
On the open torus the relative volume form is
$\Omega=dz\wedge d\log x_1\wedge d\log x_2$.
For $s>0$ the equations $y_1=\overline{x}_1$ and
$y_2=-\overline{x}_2$ cut out the vanishing sphere
$|x_1|^2+|x_2|^2=s$. The map to $z=|x_1|^2$ fibres the interior
of the interval $[0,s]$ by two-tori, with a different circle collapsing
at either end. Orient the sphere by $dz\wedge d\arg x_1\wedge d\arg x_2$.
For a holomorphic unit $g=g(z,s)$ the two angular integrations give
\begin{equation}\label{ex:nodeintegral}
 \int_{S^3}g\Omega=(2\pi i)^2\int_0^s g(z,s)\,dz.
\end{equation}
The Milnor fibre is a cotangent disc bundle of $S^3$, whose boundary
is $S^3\times S^2$ since $S^3$ is parallelizable. Inclusion of the
boundary induces an isomorphism on $H_3$. Thus
\eqref{ex:nodeintegral} evaluates the marked link class on a generator.
For $g=1+cz$ this period is $(2\pi i)^2(s+cs^2/2)$. Even on a smooth
local base, the choice of relative form can give a nonlinear period
coordinate.

\subsection*{A length-ten base}
Consider the hexagon with cyclic primitive edge vectors
\[
 (1,0),\ (1,1),\ (0,1),\ (-4,1),\ (1,-2),\ (1,-1).
\]
For edge vectors $d_i$ write $M_k=\sum_i d_i t_i^k$. For an $n$-gon
the moments through degree $n-1$ generate the full ideal. Indeed,
$M_0=0$ and the polynomial $\prod_i(T-t_i)$ gives a recurrence of
order $n$ for $M_k$. Here $t_6=0$ and $M_1=0$ give
\[
 (t_1,\ldots,t_6)=(x+5y-3z,-x-y+2z,x,y,z,0).
\]
Let $I\subset\C[x,y,z]$ be the resulting moment ideal. In graded
lexicographic order $x>y>z$, a Gr\"obner basis is
\[
\begin{gathered}
 5x^2-7xz-5y^2+5yz-z^2,\qquad
 5xy-3xz+10y^2-15yz+6z^2,\\
 3xz^2+5yz^2-z^3,\qquad y^3+yz^2-z^3,\qquad 2y^2z-z^3,\\
 yz^3,\qquad z^4.
\end{gathered}
\]
It is obtained from the two components of $M_2$ and $M_3$.
Every quartic monomial is divisible by one of its leading monomials.
By homogeneity, all moments of higher degree vanish in this quotient.
The standard monomials in degrees
two and three are $xz,y^2,yz,z^2$ and $yz^2,z^3$, respectively.
The base is supported at the origin and its Hilbert function is $(1,3,4,2)$.

On this base one can see the restriction imposed by a common change of
height. The shear $x\mapsto x+z^2$ fixing $y,z$ induces the identity on
the tangent space. For $q=x^2+xy-2xz+y^2-2yz+z^2\in I$, however,
the normal form of its pullback is
\[
 q(x+z^2,y,z)\equiv -(7yz^2+4z^3)/3\pmod I.
\]
Both monomials on the right are standard. The shear therefore fails to
restrict to the base. In the same coordinates the substitution
\[
 t_i\longmapsto H(t_i),\qquad
 H(T)=(1+x-y)T+(\alpha+z)T^2+\beta T^3,\quad \alpha,\beta\in\C,
\]
defines an automorphism. Indeed each $H(t_i)^k$ has an expansion in
$t_i$ with coefficients independent of $i$, and its moment vanishes.
The linear part of the induced endomorphism is the identity, so it is
invertible on the finite local algebra. The two edge rows give, for
the divided differences $U_i=(H(\zeta)-H(t_i))/(\zeta-t_i)$,
\[
 U_1U_2U_5U_6=U_4^4,\qquad U_2U_3U_4=U_5^2U_6.
\]
These are the identities that lift the substitution to the family.

\subsection*{The nilradical of the octagon}
For the octagon of \cite[\S9.3]{Alt} the two moment equations equate
the power sums of the triples
\[
 (t_1,t_2,t_8),\ (t_4,t_5,t_6)\qquad\text{and}\qquad
 (t_2,t_3,t_4),\ (t_6,t_7,t_8),
\]
respectively. Degrees one through three suffice by Newton's identities.
Set
\[
 (t_1,\ldots,t_8)=(a+2b+c-e,-a-b+d+e,a,b,c,d,e,0).
\]
The two linear equations now hold. Write $I\subset\C[a,b,c,d,e]$
for the four remaining equations, in degrees two and three.
The three components of the reduced base have ideals
\[
 J_1=(a-e,b),\qquad J_2=(a-d,b,c-e),\qquad J_3=(a,b-e,c).
\]
Their dimensions are $3,2,2$. To compute the thickening along their
union, put
\[
 u=ab,\qquad v=b(b-e),\qquad w=bc,\qquad
 f=(a-e)(a-d),\quad h=(a-e)(c-e).
\]
Then $J_1\cap J_2=(b,f,h)$. Modulo $J_3$, the elements $b,f,h$
become $e,ed,e^2$. It follows that
\[
 J_1\cap J_2\cap J_3=(f-db,h-eb,u,v,w).
\]
Indeed $f-db$ and $h-eb$ belong to $J_3$ and $b$ is a
nonzerodivisor modulo $J_3$. The two quadratic moment equations give
$f-db=-u-v$ and $h-eb=-2u-v-2w$ modulo $I$.

It remains to determine the powers of $(u,v,w)$ in the quotient.
Division by a Gr\"obner basis of the four moment equations, in graded
lexicographic order $a>b>c>d>e$, gives
\[
 u(u+v)=v(u+v)=uw=vw=v^3=w^3=0,\qquad
 \operatorname{NF}(v^2)=b^2d^2-b^2de-bd^2e+bde^2.
\]
Here $\operatorname{NF}$ denotes the polynomial normal form for that
order. The product relations kill every cubic in $u,v,w$. The displayed
quartic is nonzero. Since $I$ is contained in each
$J_i$, this proves that $\sqrt I/I=(u,v,w)$ and that its nilpotency
index is three. Both polygon calculations concern the local deformation
space. In a compact threefold its period image is further intersected
with the subspace of global homological relations.

\section{The first obstruction to a tangent direction}\label{sec:orders}

For $m\geq1$ put $A_m=\C[s]/(s^{m+1})$. A tangent vector extends to order $m$ if
it is induced by a deformation over $A_m$. Its first obstructed order
is the least $m$ for which no such deformation exists. For a polygon
$Q$ and $a\in W_Q$, choose the representative with $a_1=0$ and write
\[
                 M_{Q,k}(a)=\sum_i d_i a_i^k.
\]
For a local toric singularity the order-by-order criterion follows from
Altmann's equations \cite[\S2.3]{Alt}, as observed in
\cite[Appendix~A]{FHSU}. On a compact threefold the marked periods put
the local criteria in the same coordinates as the homological relations.

\begin{proposition}\label{prop:jet}
A local tangent vector $a\in W_Q$ extends to order $m$ if and only if
$M_{Q,k}(a)=0$ for $2\leq k\leq m$. Under the hypotheses of
Theorem~\ref{thm:global}, let $\xi\in T_0\Def(X)$ and let
$v=(v_p)_p\in K$ be its image under the differential of the global
period map after link duality. Set
$a_p=(\PD\circ\ell_{Q_p})^{-1}(v_p)$.
Then $\xi$ extends to order $m$ if and only if
\[
       M_{Q_p,k}(a_p)=0\qquad\text{for every $p$ and $2\leq k\leq m$.}
\]
Its first obstructed order is the least $k\geq2$ for which some
$M_{Q_p,k}(a_p)$ is nonzero. If there is no such $k$, then $\xi$ is
tangent to a one-parameter analytic deformation.
\end{proposition}

\begin{proof}
In a jet with tangent $a$, one has $t_i=s a_i+O(s^2)$. For $k\leq m$
the coefficient of $s^k$ in the $k$th equation of \eqref{eq:moments}
is $M_{Q,k}(a)$, independently of the higher coefficients of the jet.
Conversely, $t_i=s a_i$ solves all the equations over $A_m$ if the
stated moments vanish. In the product coordinates of
Theorem~\ref{thm:global}, write $\xi=(b,v)\in\C^h\oplus K$.
The same coefficient calculation applies at each $p$. For the converse,
the line $(s b,s v)$ satisfies the equations of $K$ identically, and
the remaining equations are exactly the local moment equations. If all
moments vanish, this line lies in the analytic germ itself.
\end{proof}

Petracci's pentagon \cite[Example~4.8]{Pet} has local deformation
algebra $\C[x,y]/(x^2,xy,y^3)$, so the $y$-direction first obstructs
at order three. The possible orders are unbounded even among toric
singularities whose reduced deformation space is a point.

\begin{proposition}\label{prop:orders}
For each $r\geq2$ there is an isolated Gorenstein toric threefold
singularity with Artinian semiuniversal base and a tangent vector
whose first obstructed order is $r$. The intersection of the base with
the line spanned by this vector in $W_Q$ has ring $\C[s]/(s^r)$.
\end{proposition}

\begin{proof}
For $0\leq j<r$ take the pairs of vectors
\[
        (j+2,1),\quad (-j-2,1),
\]
and adjoin $(1,1)$ and $(-1,-(2r+1))$. These primitive vectors
have distinct directions, span the plane and sum to zero. In cyclic
angular order they form the edges of a convex lattice polygon $Q$.
All but $(-1,-(2r+1))$ have vertical component one. Thus
a nonempty collection of edges with sum zero must contain that vector
and all $2r+1$ other edges. Since the edges are primitive, a
nontrivial lattice Minkowski decomposition would give a proper
nonempty collection with sum zero. Hence $B_Q$ is supported at the
origin by \cite[\S2.5]{Alt}.

Take the translation slice on which the parameter of $(-1,-(2r+1))$
is zero. Give both edges of the $j$th pair the same parameter $x_j$, and
set the parameter of $(1,1)$ to zero. On this subspace the $k$th
moment is $(0,2\sum_j x_j^k)$ for $k\geq1$. Equivalently, the
resulting section has the symmetric-group coinvariant algebra
$\C[x_0,\ldots,x_{r-1}]/(e_1,\ldots,e_r)$, where the $e_k$ are the
elementary symmetric functions, by Newton's identities. The section
is defined by equalities of edge parameters and is preserved by the
changes of volume form in Theorem~\ref{thm:local}.

Now put $x_j=s\zeta^j$, where $\zeta$ is a primitive $r$th root of
unity. The first $r-1$ moments vanish, and the $r$th is $(0,2r s^r)$.
The restricted ideal is therefore $(s^r)$. Proposition~\ref{prop:jet}
excludes every order-$r$ jet in the full base with this tangent.
\end{proof}

\section{A cubic obstruction on a projective threefold}\label{sec:compact-example}

Let $Q$ be the lattice hexagon with vertices
\[
 (0,-1),\ (1,0),\ (0,1),\ (-1,1),\ (-2,0),\ (-1,-1)
\]
in cyclic order. In $\Z^3$ take the pyramid
$P=\operatorname{conv}((Q,1),a)$ with apex $a=(0,0,-1)$.
Subdivide the face fan of $P$ at the ray through $m=(-1,0,0)$, and
let $Z$ be the resulting toric threefold.

\begin{proposition}\label{prop:compact-cubic}
A general anticanonical hypersurface $X\subset Z\times\mathbf P^1$
is a projective Calabi--Yau threefold with two singular points, both
of type $X_Q$. For $h=h^1(X,\Theta_X)$ its semiuniversal ring is
\[
 \OO_{\Def(X),0}\simeq
 \C\{w_1,\ldots,w_h,u,v,s\}/(us,vs,s^3).
\]
In particular there is a first-order deformation that extends over
$\C[\varepsilon]/(\varepsilon^3)$ and has no extension over
$\C[\varepsilon]/(\varepsilon^4)$.
\end{proposition}

\begin{proof}
The plane of the upper facet has equation $z=1$. In the order of
the edges of $Q$ the six side facets have equations
$\langle n_i,\,\cdot\,\rangle=1$ with outward normals
\[
 (2,-2,-1),\ (2,2,-1),\ (0,2,-1),\
 (-1,1,-1),\ (-1,-1,-1),\ (0,-2,-1).
\]
Thus $P$ is reflexive. The absolute determinants of the three rays
of its side cones are $1,1,1,2,2,1$, respectively. The two cones of
determinant two share the face spanned by $a$ and $(-2,0,1)$.
Its midpoint $m$ divides each into two unimodular cones. This star
subdivision is projective and crepant. The remaining maximal cone
$\sigma$ over $(Q,1)$ is smooth off its closed orbit, since the
edges of $Q$ are primitive. Write $p$ for its fixed point in $Z$.

The anticanonical bundle on $Z$ is the pullback of the globally
generated anticanonical bundle for the face fan of $P$. A general
section on $Z\times\mathbf P^1$ cuts out a smooth hypersurface away
from $\{p\}\times\mathbf P^1$ by Bertini. Its restriction there is
a general binary quadratic, with two simple roots in $\C^*$.
At each root projection to $Z$ identifies the hypersurface germ
with $(X_Q,0)$ by the implicit function theorem. In particular $X$
is normal. Adjunction gives $\omega_X\simeq\OO_X$. On the projective
Gorenstein toric fourfold $A=Z\times\mathbf P^1$ one has
$H^i(A,\OO_A)=0$ for $i>0$. Serre duality and the sequence
$0\to\omega_A\to\OO_A\to\OO_X\to0$ give
$H^0(X,\OO_X)=\C$ and $H^1(X,\OO_X)=0$.

To resolve the two points, subdivide $\sigma$ successively at
$(0,0,1)$ and $(-1,0,1)$. Their projections are the two interior
lattice points of $Q$. The first star joins the origin of the
polygon to its six edges. The second bisects the common edge of the
two triangles of normalized area two. All resulting triangles are
unimodular. Let $\widehat Z$ be the smooth projective toric
threefold so obtained and $Y\subset\widehat Z\times\mathbf P^1$
the proper transform of $X$. Near the two roots the resolution
is identified by projection with the toric resolution of $X_Q$.
Every interior edge in the triangulation contains an interior vertex
as an endpoint. Thus each exceptional fibre is the union of the two
toric divisors of the interior rays, meeting transversely along their
common invariant curve. In particular $Y\to X$ is a projective log
resolution and is an isomorphism on the regular locus.

Use the two projection charts to identify both oriented links with
the link $L$ of $X_Q$. We claim that the image of
\[
 H^2(Y,\C)\longrightarrow H^2(L,\C)\oplus H^2(L,\C)
\]
is the diagonal. Indeed the fan polytope of the original fourfold is
the convex hull of $(P,0)$ and $(0,0,0,\pm1)$. Its boundary lattice
points are the two latter vertices and the points of $(\partial P,0)$
in the lattice. The fan of $\widehat Z\times\mathbf P^1$ gives a
projective unimodular triangulation of this boundary using all its
lattice points. We may choose our general section nondegenerate with
respect to the polar Newton polytope. By the proof of
\cite[Theorem~4.4.2]{Bat}, the components of the intersections of $Y$
with the toric boundary divisors span $H^2(Y,\C)$. The six boundary
divisors corresponding to vertices of $Q$ have irreducible
intersections with $Y$, since their rays pass through vertices of the
fan polytope. They restrict in both projection charts to the same
six toric divisor classes on the punctured cone. Divisors associated
to the two interior points of $Q$ split over the two roots, and
each component is exceptional, hence disjoint from both links. All
the other boundary divisors miss the two germs. This proves
containment in the diagonal.

For the reverse inclusion, the fan of the punctured cone consists
of the proper faces of $\sigma$. Its homogeneous coordinate presentation
\cite[Theorem~2.1]{Cox} is a principal torus bundle over the punctured
cone. Indeed the fan is smooth and
its six rays span $\Z^3$, as the rays over $(1,0),(0,1),(-1,1)$
already form a lattice basis. The total space is the complement in
$\C^6$ of coordinate subspaces of complex codimension at least two
and has vanishing first and second cohomology by general position.
The low-degree sequence of the bundle then gives a surjection from
the complexified character group of the acting torus to $H^2(L,\C)$
by first Chern classes, using the retraction of the punctured cone
onto its link. The six coordinate characters generate this group
and give precisely the six ray divisor classes. This proves the claim.

By duality the relation space $K$ in \eqref{eq:K} is the
anti-diagonal in $H_2(L,\C)\oplus H_2(L,\C)$. Taking the same
marking and homogeneous toric form in the two projection charts
identifies their period images with the same germ
$G_Q=(\PD\circ\ell_Q)(B_Q)$. The restriction of a global volume form
differs from the toric generator by a unit in each chart and hence
has the same image by Theorem~\ref{thm:local}. Since the equations of $G_Q$ are
homogeneous, the pullbacks of its two ideals to $K$ coincide.
Theorem~\ref{thm:global} therefore gives $B_Q$ as the factor
transverse to the smooth germ of dimension $h$.

It remains to compute $B_Q$. The ordered edge vectors of $Q$ are
\[
 (1,1),\ (-1,1),\ (-1,0),\ (-1,-1),\ (1,-1),\ (1,0).
\]
Take the translation slice $t_6=0$. The general solution of the
first moment equation is
\[
 (t_1,\ldots,t_6)=(u+3s,v+s,4s,u+s,v+3s,0).
\]
The second moment is $(4s(u+v),4s(u-v))$. Modulo $(us,vs)$
the third is $(-12s^3,0)$. For every $k\geq 4$, setting $s=0$
cancels the two pairs with parameters $u,v$. Every other term of
the $k$th moment is divisible by $us$, $vs$ or $s^3$. This gives
the stated ring. Its quotient by $(w_1,\ldots,w_h,u,v)$ is
$\C\{s\}/(s^3)$, so $s^2\ne0$. Setting $s=\varepsilon$ and the
other parameters to zero gives a deformation modulo $\varepsilon^3$.
Every lift of its first-order class would send $s$ to
$\varepsilon+O(\varepsilon^2)$, whose cube is nonzero modulo
$\varepsilon^4$. This proves the assertion.
\end{proof}

\begin{thebibliography}{FHSU06}
\bibitem[Alt97]{Alt} K.~Altmann, \href{https://doi.org/10.1007/s002220050148}{\emph{The versal deformation of an isolated toric Gorenstein singularity}}, Invent. Math. \textbf{128} (1997), 443--479.
\bibitem[Bat94]{Bat} V.~V.~Batyrev, \href{https://arxiv.org/abs/alg-geom/9310003}{\emph{Dual polyhedra and mirror symmetry for Calabi--Yau hypersurfaces in toric varieties}}, J. Algebraic Geom. \textbf{3} (1994), 493--535.
\bibitem[Cox95]{Cox} D.~A.~Cox, \href{https://arxiv.org/abs/alg-geom/9210008}{\emph{The homogeneous coordinate ring of a toric variety}}, J. Algebraic Geom. \textbf{4} (1995), 17--50.
\bibitem[Del71]{Del} P.~Deligne, \href{https://doi.org/10.1007/BF02684692}{\emph{Th\'eorie de Hodge II}}, Publ. Math. Inst. Hautes \'Etudes Sci. \textbf{40} (1971), 5--57.
\bibitem[FHSU06]{FHSU} S.~Franco, A.~Hanany, F.~Saad and A.~M.~Uranga, \href{https://doi.org/10.1088/1126-6708/2006/01/011}{\emph{Fractional branes and dynamical supersymmetry breaking}}, J. High Energy Phys. \textbf{2006} (2006), no.~1, 011.
\bibitem[Fri86]{Fri} R.~Friedman, \emph{Simultaneous resolution of threefold double points}, Math. Ann. \textbf{274} (1986), 671--689.
\bibitem[Fri26]{Fried} R.~Friedman, \emph{Unobstructed deformations for singular Calabi--Yau varieties}, \href{https://arxiv.org/abs/2506.09857v3}{arXiv:2506.09857v3} (2026).
\bibitem[FL25]{FL} R.~Friedman and R.~Laza, \href{https://doi.org/10.4310/jdg/1755542830}{\emph{Deformations of singular Fano and Calabi--Yau varieties}}, J. Differential Geom. \textbf{131} (2025), 65--131. Numbering follows \href{https://arxiv.org/abs/2203.04823v7}{arXiv:2203.04823v7}.
\bibitem[Gro97]{Gross} M.~Gross, \href{https://doi.org/10.1007/s002080050072}{\emph{Deforming Calabi--Yau threefolds}}, Math. Ann. \textbf{308} (1997), 187--220. Numbering follows the corrected version, \href{https://arxiv.org/abs/alg-geom/9506022v2}{arXiv:alg-geom/9506022v2} (2025).
\bibitem[Gro01]{GrossSL} M.~Gross, \emph{Examples of special Lagrangian fibrations}, in \emph{Symplectic Geometry and Mirror Symmetry} (Seoul, 2000), World Scientific, 2001, 81--109, \href{https://arxiv.org/abs/math/0012002}{arXiv:math/0012002}.
\bibitem[LLW18]{LLW} Y.-P.~Lee, H.-W.~Lin and C.-L.~Wang, \emph{Towards $A+B$ theory in conifold transitions for Calabi--Yau threefolds}, J. Differential Geom. \textbf{110} (2018), 495--541, \href{https://arxiv.org/abs/1502.03277}{arXiv:1502.03277}.
\bibitem[Pet22]{Pet} A.~Petracci, \href{https://doi.org/10.1090/tran/8636}{\emph{On deformation spaces of toric singularities and on singularities of K-moduli of Fano varieties}}, Trans. Amer. Math. Soc. \textbf{375} (2022), 5617--5643.
\bibitem[Sta]{Stacks} The Stacks Project Authors, \emph{The Stacks Project}, Tags \href{https://stacks.math.columbia.edu/tag/0E4C}{0E4C} and \href{https://stacks.math.columbia.edu/tag/0C08}{0C08}.
\bibitem[Ste97]{Ste} J.~H.~M.~Steenbrink, \href{https://doi.org/10.5802/aif.1603}{\emph{Du Bois invariants of isolated complete intersection singularities}}, Ann. Inst. Fourier (Grenoble) \textbf{47} (1997), 1367--1377.
\bibitem[Wue26]{Wue} B.~J.~Wuebben, \emph{Local equations for deformations of Calabi--Yau threefolds with del Pezzo cone points}, preprint, November 26, 2026, \href{https://doi.org/10.5281/zenodo.24573015}{doi:10.5281/zenodo.24573015}.
\end{thebibliography}

\end{document}
